\chaptertitle{第二章\quad 减法——整数}

\sectiontitle{第4讲\quad 减法是加法的逆运算}

\subsection*{从一道简单的题开始}

你口袋里有5块糖。你拿出了2块给朋友。你口袋里还剩几块？

这件事用数学来写就是：$5-2=3$（读作"5减2等于3"）。

这就是\textbf{减法}：从一个数里拿走一部分，算剩下多少。

\subsection*{加法和减法是什么关系？}

现在把这3块糖放回口袋——你又有5块了。$3+2=5$。

发现没有？加法是"合在一起"，减法是"拿走"。它们是相反的两个动作：

\[
3+2=5 \quad\Longleftrightarrow\quad 5-2=3
\]

\textbf{减法是加法的逆运算。} 如果说加法是往前走，减法就是往后走。

\begin{example}
已知 $7+5=12$，不用计算直接写出下面两个结果：
$12-5=\;$? \qquad $12-7=\;$?
\end{example}
\begin{solution}
因为 $7+5=12$，把12拿走5，剩下的就是7，所以 $12-5=7$。\\
拿走7，剩下的就是5，所以 $12-7=5$。
\end{solution}

\begin{example}
已知 $25+17=42$，直接写出 $42-17$ 和 $42-25$。
\end{example}
\begin{solution}
$42-17=25$，$42-25=17$。
\end{solution}

这个关系太重要了。它告诉我们：**每一道加法题，都对应着两道减法题**。

\begin{exercise}
\item 根据加法写出两道减法：
  \begin{subexercise}
    \item $8+6=14$，$14-6=\;$? \quad $14-8=\;$?
    \item $12+9=21$，$21-9=\;$? \quad $21-12=\;$?
    \item $35+28=63$，$63-28=\;$? \quad $63-35=\;$?
    \item $123+456=579$，$579-456=\;$? \quad $579-123=\;$?
  \end{subexercise}

\item 不用计算，直接写出结果：
  \begin{subexercise}
    \item $15-8$（提示：$8+?=15$）
    \item $23-7$
    \item $50-12$
    \item $100-34$
    \item $200-56$
    \item $1000-789$
  \end{subexercise}
\end{exercise}

\subsection*{减法的竖式}

和加法一样，大数的减法也需要竖式。

\begin{example}
用竖式计算：$847-325$
\end{example}
\begin{solution}
\[
\begin{array}{r}
  847 \\
- 325 \\ \hline
\end{array}
\]

个位：$7-5=2$。十位：$4-2=2$。百位：$8-3=5$。

结果是 522。
\end{solution}

但如果某一位上面不够减呢？比如 $735-199$：

\begin{example}
用竖式计算：$735-199$
\end{example}
\begin{solution}
\[
\begin{array}{r}
  735 \\
- 199 \\ \hline
\end{array}
\]

个位：$5$减$9$不够，从十位借1个十（变成10个一），$15-9=6$。十位原来有3，被借走了1，还剩2。$2$减$9$不够，从百位借1个百（变成10个十），$12-9=3$。百位原来有7，被借走了1，还剩6。$6-1=5$。

结果是 536。
\end{solution}

这就是\textbf{借位}——减法竖式里最重要的技巧。

\subsection*{减法也有简便算法}

减法和加法一样，也有简便算法。最常用的是这条性质：

\[
a-b-c = a-(b+c)
\]

\textbf{连续减去两个数，等于减去它们的和。}

\begin{example}
计算：$1000-345-655$
\end{example}
\begin{solution}
$1000-345-655 = 1000-(345+655) = 1000-1000 = 0$。
\end{solution}

\begin{example}
计算：$735-199$
\end{example}
\begin{solution}
$199 = 200-1$，所以
$735-199 = 735-(200-1) = 735-200+1 = 535+1 = 536$。
\end{solution}

\begin{exercise}
\item 竖式计算：
  \begin{subexercise}
    \item $864-398$
    \item $1200-507$
    \item $2000-345-655$
    \item $1563-299-401$
    \item $826-79-21$
    \item $574-136-64$
  \end{subexercise}

\item 用简便方法计算：
  \begin{subexercise}
    \item $735-199$
    \item $864-398$
    \item $1685-347-253$
    \item $4321-567-433$
  \end{subexercise}

\item 应用题：
  \begin{subexercise}
    \item 一本书有 200 页，小明第一天看了 45 页，第二天看了 55 页。还剩多少页没看？（用两种方法算）
    \item 水果店有 350 个苹果，上午卖了 128 个，下午卖了 97 个。还剩多少个？
  \end{subexercise}

\item 判断：$a-b-c$ 和 $a-(b+c)$ 一定相等吗？举例说明。
\end{exercise}

\sectiontitle{第5讲\quad 负数的诞生——当减法"不够减"的时候}

\subsection*{一个让人不舒服的问题}

现在来做这道题：

\[
3-5 = \;?
\]

3块糖，要拿走5块——拿不了。3减5，不够减。

在自然数世界里，这道题没有答案。

\subsection*{但我们偏要让它有答案}

想象一根直线，你在上面标上数字。在0的位置向右走，数字变大；向左走，数字变小。

你在0的位置。$0+3$表示向右走3步，你到了3。

那 $3-5$ 呢？就是从3向左走5步——
3 $\to$ 2 $\to$ 1 $\to$ 0 $\to$ -1 $\to$ -2。

你走过0了，继续向左走，到了一个新的位置：\textbf{-2}。

\textbf{负数诞生了。} "-2"读作"负二"。

\begin{example}
用数轴找到 $1-4$ 的结果。
\end{example}
\begin{solution}
画一条数轴，标上 $-3,-2,-1,0,1,2,3,4$。从1出发向左走4步：\\
$1\to0\to-1\to-2\to-3$，所以 $1-4 = -3$。
\end{solution}

\subsection*{负数不是想象出来的}

你已经在生活中见过负数了，只是没意识到那就是负数：

\begin{itemize}
  \item 冬天温度零下5度——就是 $-5^\circ$C
  \item 你买了一个12元的东西，但你只有10元——你欠了2元，就是 $-2$元
  \item 海平面以下100米——就是 $-100$米
\end{itemize}

\subsection*{整数家族}

我们把自然数（1,2,3,\ldots）、0 和负数（-1,-2,-3,\ldots）放在一起，得到一个更大的集合，叫\textbf{整数}，用符号 $\mathbb{Z}$ 表示。

\[
\mathbb{Z} = \{\ldots,\;-3,\;-2,\;-1,\;0,\;1,\;2,\;3,\;\ldots\}
\]

\subsection*{同号两数相加}

两个正数相加我们都知道了。两个负数相加呢？

\begin{example}
计算：$(-3)+(-5)$
\end{example}
\begin{solution}
你可以这样想：欠了3元（-3），又欠了5元（-5），一共欠了8元。所以 $(-3)+(-5) = -8$。

规律：**同号两数相加，取相同的符号，把绝对值相加**。
绝对值就是"去掉负号后的数"。
\end{solution}

\begin{example}
计算：$(-23)+(-15)$
\end{example}
\begin{solution}
$23+15=38$，都是负数，结果是 $-38$。
\end{solution}

\subsection*{异号两数相加}

一个正数加一个负数呢？

\begin{example}
计算：$(-5)+3$
\end{example}
\begin{solution}
你欠了5元（-5），还了3元（+3），现在你还欠2元。所以 $(-5)+3 = -2$。

规律：**异号两数相加，取绝对值较大的符号，用大的绝对值减去小的绝对值**。
$|-5|=5$，$|3|=3$，$5>3$，取负号，$5-3=2$，结果是 $-2$。
\end{solution}

\begin{example}
计算：$5+(-3)$
\end{example}
\begin{solution}
你有5元（+5），花了3元（-3），还剩2元。$5+(-3)=2$。
取正号，$5-3=2$。
\end{solution}

\begin{exercise}
\item 用数轴找出结果：
  \begin{subexercise}
    \item $2-7$
    \item $1-5$
    \item $0-4$
    \item $-1-3$
  \end{subexercise}

\item 计算：
  \begin{subexercise}
    \item $(-3)+(-7)$
    \item $(-15)+(-23)$
    \item $(-5)+8$
    \item $12+(-7)$
    \item $(-3.2)+4.5$
    \item $(-2.5)+(-1.5)$
  \end{subexercise}

\item 温度问题：
  \begin{subexercise}
    \item 早晨温度是 $5^\circ$C，中午升了 $8^\circ$C，晚上降了 $15^\circ$C。晚上的温度是多少？
    \item 冰箱冷冻室的温度是 $-18^\circ$C，调到冷藏模式后升高了 $25^\circ$C。现在的温度是多少？
  \end{subexercise}

\item 金钱问题：
  \begin{subexercise}
    \item 小明有10元，买了一支12元的笔。他现在有多少钱？
    \item 小红欠了小刚5元，她又向小刚借了3元。小红一共欠小刚多少元？
  \end{subexercise}

\item 判断：
  \begin{subexercise}
    \item 两个自然数相减，结果一定是自然数吗？
    \item 两个整数相加，结果一定是整数吗？
    \item 两个整数相减，结果一定是整数吗？
  \end{subexercise}
\end{exercise}

\sectiontitle{第6讲\quad 绝对值和整数加减混合运算}

\subsection*{绝对值——只看距离，不看方向}

看看下面两个数：3 和 -3。

在数轴上，3 在0的右边3格处，-3 在0的左边3格处。它们的方向相反，但有一点相同：\textbf{它们离0一样远}。

这个"离0有多远"的数，叫做\textbf{绝对值}。用 $|a|$ 表示。

\[
|3| = 3,\qquad |-3| = 3,\qquad |0| = 0
\]

绝对值的正式定义：一个数在数轴上的点到原点的距离。

\begin{example}
计算：$|-5|+|3|$
\end{example}
\begin{solution}
$|-5|=5$，$|3|=3$，$5+3=8$。
\end{solution}

\begin{example}
比较大小：$-3$ 和 $-5$
\end{example}
\begin{solution}
在数轴上，-3 在 -5 的右边。数轴上右边的数总比左边的大，所以 $-3 > -5$。

注意：在负数中，看起来"更大"的数（比如-3）其实绝对值更小。这是一个容易搞混的地方。
\end{solution}

\subsection*{整数加减混合运算}

现在我们有了正数、负数和0。混合运算时，有一条法则：

\textbf{减去一个数，等于加上这个数的相反数}。

\[
a - b = a + (-b)
\]

这条法则让我们可以把所有减法都变成加法（加法比减法好处理）：

\begin{example}
计算：$28-(-14)+(-18)$
\end{example}
\begin{solution}
先处理符号：减去 $-14$ 等于加上 $14$。所以
$28-(-14)+(-18) = 28+14+(-18)$。
然后同号先加：$28+14=42$，$42+(-18)=24$。
结果是24。
\end{solution}

\begin{example}
计算：$(-12)-(-15)+(-8)-(-20)$
\end{example}
\begin{solution}
依次转化：$(-12)+15+(-8)+20$。
同号先加：$(-12)+(-8) = -20$，$15+20=35$。
$(-20)+35 = 15$。
\end{solution}

\begin{example}
计算：$(-15)-(-23)+(-18)-(-7)$
\end{example}
\begin{solution}
依次转化：$(-15)+23+(-18)+7$。
同号先加：$(-15)+(-18) = -33$，$23+7=30$。
$(-33)+30 = -3$。
\end{solution}

\begin{example}
计算：$35-(-28)+(-42)-(-15)$
\end{example}
\begin{solution}
转化：$35+28+(-42)+15$。
同号先加：$35+28+15 = 78$，$78+(-42)=36$。
结果是36。
\end{solution}

\subsection*{绝对值比大小的几个例子}

负数之间比大小要特别小心。数轴上越靠左的数越小——所以 $-10 < -1$，虽然10看起来比1大。

\begin{example}
比较 $-3$ 和 $-7$ 的大小。
\end{example}
\begin{solution}
在数轴上，$-3$ 在 $-7$ 的右边（离0更近）。右边总比左边大，所以 $-3 > -7$。
\end{solution}

\begin{example}
把 $-5,\;2,\;0,\;-1,\;4$ 按从小到大的顺序排列。
\end{example}
\begin{solution}
先找负数：$-5$ 和 $-1$。$-5 < -1 < 0$。再加正数：$0 < 2 < 4$。
完整的顺序：$-5 < -1 < 0 < 2 < 4$。
\end{solution}

\subsection*{运算律在整数中仍然成立}

加法交换律、结合律在整数中依然成立。凑整法也继续有效——只是现在要注意符号。

\begin{example}
用凑整法计算：$(-19)+(-35)+(-21)$
\end{example}
\begin{solution}
把能凑整的放一起：$(-19)+(-21) = -40$，$-40+(-35) = -75$。
\end{solution}

\begin{example}
用凑整法计算：$87+(-46)+(-54)+13$
\end{example}
\begin{solution}
$87+13=100$，$(-46)+(-54)=-100$，$100+(-100)=0$。
\end{solution}

\begin{exercise}
\item 计算绝对值：
  \begin{subexercise}
    \item $|-5|$
    \item $|-8|$
    \item $|0|$
    \item $|3-5|$
    \item $|-5|+|3|$
    \item $|-5+3|$
  \end{subexercise}

\item 比较大小：
  \begin{subexercise}
    \item $-7$ 和 $-3$
    \item $-2$ 和 $0$
    \item $-|{-3}|$ 和 $-(-3)$
    \item $-15$ 和 $-5$
    \item $-|{-8}|$ 和 $-(-8)$
    \item $-(-4)$ 和 $-(+2)$
  \end{subexercise}

\item 从小到大排列：
  \begin{subexercise}
    \item $-3,\;5,\;0,\;-1,\;2$
    \item $-7,\;4,\;-2,\;1,\;-5$
    \item $0,\;-8,\;3,\;-3,\;6$
  \end{subexercise}

\item 计算：
  \begin{subexercise}
    \item $(-15)+(-23)+34$
    \item $28-(-14)+(-18)$
    \item $(-12)-(-15)+(-8)-(-20)$
    \item $(-18)-(-22)+(-5)-(-8)$
    \item $(-30)+(-25)+(+45)+(-18)$
    \item $12+(-35)+(-18)+(+25)$
  \end{subexercise}

\item 含绝对值的计算：
  \begin{subexercise}
    \item $|{-5}|+|{-3}|$
    \item $|{-8}|-|{-2}|$
    \item $|{-3}|\times|{-4}|$
    \item $|{-12}|\div|{-3}|$
    \item $|-2^2|+|-3^2|$
    \item $|-5|-|{-3}^2|$
  \end{subexercise}

\item 已知条件题：
  \begin{subexercise}
    \item 已知 $|x|=2$，$|y|=3$，求 $x+y$ 的值（注意多种可能）。
    \item 已知 $|x-2|+|y+3|=0$，求 $x$、$y$ 的值。
  \end{subexercise}

\item 若 $a$、$b$ 互为相反数，$c$、$d$ 互为倒数，$|m|=2$，求 $\dfrac{a+b}{m}+m^2-cd$ 的值（提示：$a+b=0$，$cd=1$，$m^2=4$）。
\end{exercise}

\sectiontitle{章末小结}

这一章我们经历了数学史上最重要的一次"扩张"：

\begin{enumerate}
  \item \textbf{减法是加法的逆运算}——合起来和拆回去的关系。
  \item \textbf{减法不够减？负数来了。} 3-5 不是"算不出"，而是 $-2$。
  \item \textbf{整数 $\mathbb{Z}$} = 自然数 + 0 + 负数。
  \item \textbf{绝对值}是一个数到0的距离，不看方向。
  \item \textbf{整数加减混合运算}的法则：同号相加取相同符号，异号相加取绝对值大的符号；减法全都可以转化为加法。
\end{enumerate}

\textbf{关键思想：}运算会"不够用"。不够用的时候，不要硬算，而是扩张数的范围让它变得可以算。这个思路后面还会出现两次。

\bigskip

\begin{center}
\fbox{
  \parbox{0.7\linewidth}{
    \textbf{数系脚印 \#2：} $\mathbb{N} \to \mathbb{Z}$ \\
    因为减法不够用，我们加上了负数。\\
    但还有更多"不够用"在前面等着。
  }
}
\end{center}
